\ifdefined\IntervalProjectMaster
\def\IntervalProjectEnd{}
\section{The open two-interval endpoint case: the \texorpdfstring{$R_5$}{R5}-first proof}
\label{sec:open-two-interval-r5}
\else
\documentclass{article}
\usepackage[utf8]{inputenc}
\usepackage{amsmath}
\usepackage{amssymb}
\usepackage{amsthm}
\usepackage{parskip}

\newtheorem{theorem}{Theorem}
\newtheorem{claim}{Claim}

\title{The endpoint case for two open intervals}
\author{Ferran Espu\~na}
\date{\today}
\def\IntervalProjectEnd{\end{document}}

\begin{document}

\maketitle
\fi

\begin{theorem}\label{thm:open_two_intervals_endpoint}
Let $A \subset \mathbb T$ be a union of two non-empty disjoint open intervals with
\[
    \mu(A)=\frac15.
\]
If
\[
    A+A-3\cdot A \neq \mathbb T,
\]
then, up to translating $A$ and replacing $A$ by $-A$, we have
\[
    A=\left(0,\frac1{10}\right)\cup\left(\frac12,\frac35\right).
\]
In particular, $2\cdot A$ is contained in an interval of length $\frac15$.
\end{theorem}

\begin{proof}
The property of $A+A-3\cdot A$ being equal to $\mathbb T$ is unchanged by translating
$A$ or by replacing $A$ by $-A$. As in the closed-interval proof, we may therefore write
\[
    A=I\cup J, \qquad I=(0,\alpha), \qquad J=(x,x+\beta),
\]
where
\begin{align}
    0<\alpha < x < x+\beta <1, \label{cond:order} \\
    \alpha \geq \beta, \label{cond:alpha_geq_beta_open_R5} \\
    x-\alpha \leq 1-(x+\beta). \label{cond:shorter_gap}
\end{align}
Since $\mu(A)=\frac15$, we have
\[
    \alpha+\beta=\frac15.
\]
Condition~\eqref{cond:shorter_gap} is equivalent to
\begin{equation}\label{ineq:x_leq_3a_2b}
    x \leq \frac{1+\alpha-\beta}{2}=3\alpha+2\beta.
\end{equation}
We will show that, unless
\begin{equation}\label{eq:exceptional_case}
    \alpha=\beta=\frac1{10}, \qquad x=\frac12,
\end{equation}
we have $A+A-3\cdot A=\mathbb T$.

Let
\begin{align*}
    R_1=I+I-3\cdot I &= (-3\alpha,2\alpha), \\
    R_2=I+J-3\cdot I &= (x-3\alpha,x+\alpha+\beta), \\
    R_3=J+J-3\cdot I &= (2x-3\alpha,2x+2\beta), \\
    R_4=I+I-3\cdot J &= (-3x-3\beta,-3x+2\alpha), \\
    R_5=I+J-3\cdot J &= (-2x-3\beta,\alpha+\beta-2x), \\
    R_6=J+J-3\cdot J &= (-x-3\beta,-x+2\beta).
\end{align*}
All of these intervals have length strictly smaller than $1$, so we use them in proper
form: a point $y\in\mathbb T$ belongs to one of these intervals if and only if, for
some $k\in\mathbb Z$, the real number $y+k$ belongs to the displayed interval.
Clearly each $R_i$ is contained in $A+A-3\cdot A$.

We first record the strict inequalities that replace the non-strict endpoint inequalities
in the closed-interval proof.

\begin{claim}\label{claim:strict_first_inequalities}
If~\eqref{eq:exceptional_case} does not hold, then
\[
    x<5\alpha
    \qquad\text{and}\qquad
    x<4\alpha+\beta.
\]
\end{claim}

\begin{proof}
By~\eqref{ineq:x_leq_3a_2b} and $\alpha\geq\beta$,
\[
    x\leq 3\alpha+2\beta\leq 5\alpha.
\]
If equality held in $x\leq 5\alpha$, then we would need $\alpha=\beta$ and
$x=3\alpha+2\beta$. Since $\alpha+\beta=\frac15$, this gives
$\alpha=\beta=\frac1{10}$ and $x=\frac12$, which is exactly
\eqref{eq:exceptional_case}. Hence, outside the exceptional case, $x<5\alpha$.

Similarly,
\[
    x\leq 3\alpha+2\beta\leq 4\alpha+\beta,
\]
and equality again forces $\alpha=\beta=\frac1{10}$ and $x=\frac12$. Thus
$x<4\alpha+\beta$ outside the exceptional case.
\end{proof}

\begin{claim}\label{claim:first_cover}
Assume~\eqref{eq:exceptional_case} does not hold. Then
\[
    R_1\cup R_2\cup R_3
\]
covers the interval
\[
    [0,2x+2\beta)
\]
and also covers the interval $(1-3\alpha,1)$.
\end{claim}

\begin{proof}
Modulo $1$, the interval $R_1=(-3\alpha,2\alpha)$ covers
\[
    [0,2\alpha)\cup(1-3\alpha,1).
\]
By Claim~\ref{claim:strict_first_inequalities}, $x<5\alpha$, so
\[
    x-3\alpha<2\alpha.
\]
Also $2\alpha<x+\alpha+\beta$, since $\alpha<x+\beta$. Hence $R_2$ contains
$2\alpha$, and therefore $R_1\cup R_2$ covers $[0,x+\alpha+\beta)$.

Again by Claim~\ref{claim:strict_first_inequalities}, $x<4\alpha+\beta$, so
\[
    2x-3\alpha<x+\alpha+\beta.
\]
Also $x+\alpha+\beta<2x+2\beta$, since $\alpha<x+\beta$. Hence $R_3$ contains
$x+\alpha+\beta$, and $R_1\cup R_2\cup R_3$ covers $[0,2x+2\beta)$.
\end{proof}

It remains to cover the possible leftover interval
\begin{equation}\label{eq:leftover}
    [2x+2\beta,1-3\alpha].
\end{equation}
If this interval is empty, then we are done. Thus assume from now on that it is non-empty.
Then
\begin{equation}\label{ineq:leftover_nonempty}
    2x+2\beta\leq 1-3\alpha,
\end{equation}
or equivalently, using $\alpha+\beta=\frac15$,
\begin{equation}\label{ineq:two_x_bound}
    2x\leq 2\alpha+3\beta.
\end{equation}

We now follow the original proof and try to cover the leftover using $R_4\cup R_5$.
The only difference is that the intervals are open, so one endpoint case survives.

\begin{claim}\label{claim:R4_R5_cover_except_fake}
The intervals $R_4$ and $R_5$ cover the leftover interval~\eqref{eq:leftover}, except
possibly in the case
\begin{equation}\label{eq:fake_case_R5}
    \alpha=\beta=\frac1{10}, \qquad x=\frac14.
\end{equation}
In that case, the only point of~\eqref{eq:leftover} not covered by $R_4\cup R_5$ is
$\frac7{10}$.
\end{claim}

\begin{proof}
We check the leftover interval after shifting by $-1$. Thus we want to cover
\[
    [2x+2\beta-1,-3\alpha].
\]
The intervals $R_4$ and $R_5$ overlap, because
\[
    -2x-3\beta<-3x+2\alpha
\]
is equivalent to $x<2\alpha+3\beta$, which follows from
\eqref{ineq:two_x_bound}. Therefore
\[
    R_4\cup R_5=(-3x-3\beta,\alpha+\beta-2x).
\]
The left endpoint is strictly covered, since
\[
    -3x-3\beta<2x+2\beta-1
\]
is equivalent to $x+\beta>\frac15$, which follows from $x>\alpha$ and
$\alpha+\beta=\frac15$.

For the right endpoint, we need
\[
    -3\alpha<\alpha+\beta-2x,
\]
which is equivalent to
\[
    2x<4\alpha+\beta.
\]
By~\eqref{ineq:two_x_bound} and $\alpha\geq\beta$,
\[
    2x\leq 2\alpha+3\beta\leq 4\alpha+\beta.
\]
Thus the desired strict inequality can fail only if both inequalities are equalities.
This forces $\alpha=\beta=\frac1{10}$ and $2x=2\alpha+3\beta=\frac12$, i.e.
$x=\frac14$.

In that case, the leftover interval~\eqref{eq:leftover} is the singleton
\[
    \left\{\frac7{10}\right\},
\]
and $\frac7{10}-1=-\frac3{10}$ is exactly the right endpoint of $R_5$, so it is not
covered by $R_4\cup R_5$.
\end{proof}

This is the only place where the patch by $R_6$ is needed. In the fake case
\eqref{eq:fake_case_R5}, we have
\[
    R_6=\left(-\frac{11}{20},-\frac1{20}\right),
\]
and
\[
    \frac7{10}-1=-\frac3{10}\in R_6.
\]
Therefore the apparent missing endpoint is covered by $R_6=J+J-3\cdot J$.

We have shown that, outside the exceptional case~\eqref{eq:exceptional_case}, the union
of the intervals $R_1,\ldots,R_6$ covers all of $\mathbb T$. Hence
$A+A-3\cdot A=\mathbb T$ outside~\eqref{eq:exceptional_case}. Therefore, if
$A+A-3\cdot A\neq\mathbb T$, the exceptional case~\eqref{eq:exceptional_case} must hold.

Finally, in the exceptional case,
\[
    A=\left(0,\frac1{10}\right)\cup\left(\frac12,\frac35\right),
\]
and hence
\[
    2\cdot A=\left(0,\frac15\right),
\]
which is contained in an interval of length $\frac15$. A direct calculation also shows
that this exceptional set really does not cover the circle: the six intervals above have
union
\[
    \mathbb T\setminus\left\{\frac15,\frac7{10}\right\}.
\]
This completes the proof.
\end{proof}

\IntervalProjectEnd
